% Иллюстрация контрастивного обучения и формирования пакета (Рисунок 4.2 диссертации).
% Собирается отдельно из этой папки: pdflatex figure_3.tex -> figure_3.pdf
\documentclass[tikz,border=4pt]{standalone}
\usepackage[T2A]{fontenc}
\usepackage[utf8]{inputenc}
\usepackage[russian]{babel}
\usepackage{amsmath}
\usepackage{fontawesome5}
\usetikzlibrary{arrows.meta,calc}

% заливки блоков (как в figure_1.tex)
\definecolor{yel}{RGB}{255,242,204}  \definecolor{blu}{RGB}{218,232,252}
\definecolor{grn}{RGB}{213,232,212}  \definecolor{pur}{RGB}{225,213,231}
\definecolor{gry}{RGB}{226,226,226}  \definecolor{card}{RGB}{255,247,224}
\definecolor{edge}{RGB}{120,120,120} \definecolor{conn}{RGB}{95,95,95}
\definecolor{fire}{RGB}{232,96,28}   \definecolor{ice}{RGB}{66,146,214}
\definecolor{wave}{RGB}{60,130,70}
\definecolor{cellA}{RGB}{204,220,246} \definecolor{cellB}{RGB}{152,182,230} \definecolor{cellC}{RGB}{100,140,205}
\definecolor{diag}{RGB}{214,228,250}

% блок-энкодер (сужается вправо): x0 x1 ytop ybot заливка
\newcommand{\encblock}[5]{\draw[blk, fill=#5] (#1,#3) -- (#2,#3-0.45) -- (#2,#4+0.45) -- (#1,#4) -- cycle;}
% стопка отсчетов стимула: x y нижнего левого угла передней карточки (1.6 x 0.9)
\newcommand{\cardstack}[2]{\foreach \i in {2,1,0} {\fill[card, draw=edge, line width=0.5pt, rounded corners=1pt] ($(#1,#2)+(\i*0.15,\i*0.15)$) rectangle ($(#1,#2)+(\i*0.15+1.6,\i*0.15+0.9)$);}}
% модальность 1, пространственный сигнал: стопка тензоров с картой 3x3 (1.5 x 1.0)
\newcommand{\tensoricon}[2]{%
  \foreach \i in {2,1} {\fill[blu, draw=edge, line width=0.5pt] ($(#1,#2)+(\i*0.15,\i*0.15)$) rectangle ($(#1,#2)+(\i*0.15+1.5,\i*0.15+1.0)$);}
  \foreach \c/\r/\col in {0/0/cellB,1/0/cellA,2/0/cellC,0/1/cellA,1/1/cellC,2/1/cellB,0/2/cellC,1/2/cellB,2/2/cellA} {
    \fill[\col] (#1+\c*0.5,#2+\r*0.3333) rectangle (#1+\c*0.5+0.5,#2+\r*0.3333+0.3333);}
  \foreach \k in {1,2} {\draw[white, line width=0.6pt] (#1+\k*0.5,#2) -- (#1+\k*0.5,#2+1.0) (#1,#2+\k*0.3333) -- (#1+1.5,#2+\k*0.3333);}
  \draw[edge, line width=0.5pt] (#1,#2) rectangle (#1+1.5,#2+1.0);}
% модальность 2, многоканальный временной ряд: окно с четырьмя гладкими кривыми (1.7 x 1.2)
\newcommand{\seriesicon}[2]{%
  \fill[white, draw=edge, line width=0.5pt] (#1,#2) rectangle (#1+1.7,#2+1.2);
  \begin{scope}\clip (#1+0.04,#2+0.04) rectangle (#1+1.66,#2+1.16);
    \foreach \k/\ph in {0/0,1/110,2/200,3/300} {
      \draw[wave, line width=0.55pt] plot[domain=0:1.5, samples=40, smooth] ({#1+0.1+\x},{#2+1.02-\k*0.28+0.09*sin(\x*330+\ph)});}
  \end{scope}}

\begin{document}
\begin{tikzpicture}[font=\small, line cap=round, line join=round,
  blk/.style={draw=edge, line width=0.55pt},
  cell/.style={draw=edge, line width=0.45pt},
  lbl/.style={align=center, inner sep=1pt},
  hdr/.style={font=\normalsize\bfseries, align=center, inner sep=1pt},
  con/.style={conn, line width=0.6pt},
  arr/.style={con, ->, >={Stealth[length=3pt,width=2.6pt]}},
  ic/.style={font=\normalsize, inner sep=0pt},
]

% ---------- легенда и заголовок ----------
\node[ic, text=fire] at (0.8,-0.6) {\faFire};
\node[lbl, anchor=west] at (1.1,-0.6) {Обучаемый};
\node[ic, text=ice] at (3.5,-0.6) {\faSnowflake};
\node[lbl, anchor=west] at (3.8,-0.6) {Замороженный};
\node[hdr] at (11.2,-0.6) {Контрастивное обучение: пакет из $B$ триплетов};

% ---------- стимул и его энкодер ----------
\cardstack{0.6}{-2.85}
\node[lbl] at (1.55,-3.2) {Стимул};
\encblock{2.9}{4.9}{-1.3}{-3.63}{yel}
\node[lbl] at (3.9,-2.47) {Энкодер\\стимула};
\node[ic, text=ice] at (4.75,-1.65) {\faSnowflake};
\draw[con] (2.55,-2.47) -- (2.9,-2.47);

% шина к строке представлений стимула
\draw[con] (4.9,-2.47) -- (15.125,-2.47);
\foreach \x in {10.125,11.375,12.625,15.125} {\draw[arr] (\x,-2.47) -- (\x,-3.25);}
\node[lbl] at (13.875,-2.85) {$\cdots$};
\foreach \j/\t in {0/1,1/2,2/3,4/B} {
  \draw[cell, fill=yel] (9.5+\j*1.25,-4.1) rectangle (10.75+\j*1.25,-3.3);
  \node at (10.125+\j*1.25,-3.7) {$\mathbf{z}_I^{\t}$};}
\draw[cell, fill=yel] (13.25,-4.1) rectangle (14.5,-3.3);
\node at (13.875,-3.7) {$\cdots$};
\node[lbl, anchor=west] at (15.95,-3.7) {Представления\\стимула};

% ---------- модальности, энкодеры, модуль объединения ----------
\tensoricon{0.65}{-5.6}
\node[lbl] at (1.55,-5.98) {Модальность 1};
\encblock{2.9}{4.9}{-3.85}{-6.18}{blu}
\node[lbl] at (3.9,-5.02) {Энкодер 1};
\node[ic, text=fire] at (4.75,-4.2) {\faFire};
\draw[con] (2.5,-5.02) -- (2.9,-5.02);

\seriesicon{0.7}{-8.1}
\node[lbl] at (1.55,-8.48) {Модальность 2};
\encblock{2.9}{4.9}{-6.35}{-8.68}{grn}
\node[lbl] at (3.9,-7.52) {Энкодер 2};
\node[ic, text=fire] at (4.75,-6.7) {\faFire};
\draw[con] (2.45,-7.5) -- (2.9,-7.5);

\encblock{5.6}{7.6}{-5.1}{-7.43}{pur}
\node[lbl] at (6.6,-6.27) {Модуль\\объединения};
\node[ic, text=fire] at (7.45,-5.45) {\faFire};
\draw[con] (4.9,-5.02) to[out=0,in=180] (5.6,-5.95);
\draw[con] (4.9,-7.52) to[out=0,in=180] (5.6,-6.6);

% ---------- столбец объединенных представлений ----------
\draw[con] (7.6,-6.27) -- (8.0,-6.27);
\draw[con] (8.0,-5.05) -- (8.0,-8.65);
\foreach \i in {0,1,2,4} {\draw[arr] (8.0,-5.05-\i*0.9) -- (8.55,-5.05-\i*0.9);}
\foreach \i/\t in {0/1,1/2,2/3,4/B} {
  \draw[cell, fill=pur] (8.6,-4.6-\i*0.9) rectangle (9.3,-5.5-\i*0.9);
  \node at (8.95,-5.05-\i*0.9) {$\mathbf{z}_c^{\t}$};}
\draw[cell, fill=pur] (8.6,-7.3) rectangle (9.3,-8.2);
\node at (8.95,-7.75) {$\vdots$};
\node[lbl, anchor=north] at (8.95,-9.25) {Объединенные\\представления};

% ---------- матрица логитов ----------
\foreach \i in {0,...,4} {\foreach \j in {0,...,4} {
  \draw[cell, fill=white] (9.5+\j*1.25,-4.6-\i*0.9) rectangle (10.75+\j*1.25,-5.5-\i*0.9);}}
\foreach \k in {0,...,4} {\draw[cell, fill=diag] (9.5+\k*1.25,-4.6-\k*0.9) rectangle (10.75+\k*1.25,-5.5-\k*0.9);}
\draw[cell, fill=gry] (9.5,-8.2) rectangle (10.75,-9.1);
\foreach \i/\a in {0/1,1/2,2/3,4/B} {\foreach \j/\b in {0/1,1/2,2/3,4/B} {
  \node at (10.125+\j*1.25,-5.05-\i*0.9) {$\mathbf{z}_c^{\a}\!\cdot\mathbf{z}_I^{\b}$};}}
\foreach \j in {0,1,2,4} {\node at (10.125+\j*1.25,-7.75) {$\vdots$};}
\foreach \i in {0,1,2,4} {\node at (13.875,-5.05-\i*0.9) {$\cdots$};}
\node at (13.875,-7.75) {$\ddots$};
\draw[cell, fill=diag] (16.0,-5.25) rectangle (16.4,-4.85);
\node[lbl, anchor=west] at (16.5,-5.05) {Положительные\\пары};
\draw[cell, fill=white] (16.0,-6.6) rectangle (16.4,-6.2);
\node[lbl, anchor=west] at (16.5,-6.4) {Отрицательные\\пары};

% ---------- формула для выделенной пары ----------
\draw[con, densely dashed, ->, >={Stealth[length=3pt,width=2.6pt]}] (10.5,-9.1) -- (10.5,-9.9) -- (11.55,-9.9);
\node[lbl, anchor=west] at (11.65,-9.9) {$\dfrac{\langle \mathbf{z}_c^{B},\, \mathbf{z}_I^{1}\rangle}{\|\mathbf{z}_c^{B}\|\,\|\mathbf{z}_I^{1}\|}\; e^{\tau}$};

\end{tikzpicture}
\end{document}
